%% Cover letter for the International Journal of Information Security (Springer).
%% Compile with pdfLaTeX. Red [MISSING] items must be completed by the authors.
\documentclass[10pt]{letter}
\usepackage[a4paper,margin=2cm]{geometry}
\usepackage{xcolor}
\usepackage[hidelinks]{hyperref}
\newcommand{\MISSING}[1]{\textcolor{red}{[MISSING: #1]}}
\setlength{\parskip}{0.35em}

\signature{Mounish Senisetty, on behalf of all authors\\ \MISSING{affiliation}\\ senisettymounish@gmail.com}
\address{\MISSING{Department, Institution}\\ \MISSING{City, Country}}
\date{\today}

\begin{document}
\begin{letter}{The Editor-in-Chief\\ \emph{International Journal of Information Security}\\ Springer Nature}

\opening{Dear Editor-in-Chief,}

We submit ``Bounded-memory streaming redaction of personal data and credentials with a whole-document equivalence guarantee'' for consideration as an original research article in the \emph{International Journal of Information Security}.

Logs, database exports and packet captures must be stripped of personal data and credentials before they are shared, and they are routinely too large to process in one pass. Chunked processing introduces a silent leak: a secret that straddles a chunk boundary is missed and written out in the clear. We found no redaction, de-identification or secret-scanning system that states when its chunked output equals its whole-document output.

The paper closes that gap for context-dependent redaction. Each detector declares the longest span it emits and the context it reads; we prove that an overlap covering every detector's span plus its right context makes streamed output equal to whole-document output, that the parallel engine is byte-identical to the sequential one or falls back to it, and that working memory is independent of input size. All proofs are in the paper (Section~5 and Appendix~A), each step mapped to the implementation. The default overlap meets the requirement exactly. Measured memory stays between 8.4 and 8.54\,MiB from 8\,MiB to 2\,GiB of input, against 904.6\,MiB for whole-document redaction of a 32\,MiB file.

We have tried to report the evaluation as a referee would want it. Detection quality is measured on held-out splits against Presidio (three configurations), gitleaks and TruffleHog, with paired bootstrap intervals computed at the document level and, to account for dependence, with cluster and two-stage resampling by rule, template and domain. We report where margins do not survive clustering, the weak document-level protection on RedactionBench, and a contamination analysis that compares pre- and post-exposure code. Code, split manifests, raw results and per-document predictions are public.

The work fits the journal's scope in data protection, privacy-enhancing technologies and secure data handling, and should interest readers who build or deploy data-loss-prevention and secret-scanning pipelines.

\MISSING{Keep this paragraph only if the FTNCT'09 paper was submitted or published; complete it and attach a copy, or delete it.} A preliminary version of the streaming architecture appeared as \MISSING{full citation}. This article adds the per-detector formal analysis with full proofs, held-out evaluation, secret-scanner baselines, cluster-robust statistics, an ablation, a contamination analysis and the memory curve.

The manuscript is not under consideration elsewhere, and all authors approve the submission. As stated in the manuscript, an AI assistant was used to implement evaluation scripts and to draft text; the authors reviewed every result and every sentence and take full responsibility for the content. We declare no competing interests \MISSING{confirm}.

\closing{Yours sincerely,}

\end{letter}
\end{document}
